\section{统一比较矩阵}

第 6、7 章分开讲了“继承什么”与“借入什么”。本章把这条研究线上的主要方法放进一张矩阵横向比较：发刷、多项式 Wolff 公理、平面刷、粘性纲领、三维全解、OpenAI 三维篇与四维篇。比较的四个轴是目标与维度、结论强度、假设、核心工具，另附跨门派成分与验证状态两列。

\begin{table}[t]
\centering
\caption{挂谷研究主要方法的统一比较矩阵}
\label{tab:matrix}
\scriptsize
\begin{tabularx}{\textwidth}{@{}L{2.0cm}L{2.3cm}L{2.5cm}XL{1.5cm}L{1.7cm}@{}}
\toprule
\rowcolor{lightbg}\textbf{方法（年份）} & \textbf{目标} & \textbf{结论强度} & \textbf{核心工具} & \textbf{跨门派成分} & \textbf{验证状态} \\
\midrule
发刷（1995）\upcite{wolff1995} & $\mathbb{R}^n$ 维数下界 & $(n+2)/2$；$\mathbb{R}^4$ 得 $3$ & 发刷、管交体积估计 & 低 & 已证、同行评审 \\
多项式 Wolff 公理（2018）\upcite{guthzahl2018,zahl2018,katzrogers2018} & $\mathbb{R}^4$ 维数下界 & $3+\frac{1}{40}$ & 多项式划分、Wolff 公理、Severi 型代数几何 & 低（有限域方法早已内化） & 已证、同行评审 \\
平面刷（2021）\upcite{katz2021} & $\mathbb{R}^4$ 维数与极大参数 & $3.059$ / $3.049$ & 平面刷、阴影计数、多尺度组织 & 低 & 已证、同行评审 \\
粘性纲领（2022--2026）\upcite{wangz2022,raichoudhuri2026,wangzakharov2026} & 粘性集的维数 & $\mathbb{R}^3$ 全维（限粘性）；$\mathbb{R}^4$ 粘性 $13/4$ & 多尺度归纳、加权计数、稠密阴影条件 & 低 & 已证（条件性结论） \\
三维全解（2025--2026）\upcite{wangz2025,gwz2026} & $\mathbb{R}^3$ 集合 & $=3$（Hausdorff 与 Minkowski） & 凸聚类、棱柱$\to$板条、阴影、加权并集估计 & 低 & 已证、同行评审（流线化版本 2026） \\
OpenAI 三维篇（2026）\upcite{openai3d2026} & $\mathbb{R}^3$ 极大函数 & $\lambda^{3}$（临界密度幂）；$L^3\to L^3$ 带 $\delta^{-\varepsilon}$ & 加权全时段板条 + 熵账本 + 耦合 + 钉住 + Furstenberg & 中高 & 未形式化、未评审（声称） \\
OpenAI 四维篇（2026）\upcite{openai4d2026} & $\mathbb{R}^4$ Hausdorff 维数 & $=4$（任意集合） & 二次图卡 + 地平线 + 标量投影 + 借入工具集 & 高 & 未形式化、未评审（声称） \\
\bottomrule
\end{tabularx}
\end{table}

\subsection{三条读表线索}

\textbf{线索一：假设在收缩。}从粘性纲领到三维全解，条件从“粘性、稠密阴影”收缩到无条件；四维篇承接同样姿态，明确不设可测性、紧性与粘性假设\upcite{openai4d2026}。假设收缩意味着结论覆盖面变宽，也意味着证明必须处理更野的配置——这是四维篇体量膨胀的直接原因。

\textbf{线索二：结论在精确化。}三维篇把遗留的密度幂 $\lambda^{K(\varepsilon)}$ 精化为 $\lambda^{3}$，属于“在已知框架内把最优点补圆”；四维篇把下界从 $3.059$ 直接声称到 $4$，属于“把量级从改进变为补全”。矩阵里这两格的差异，对应两种性质完全不同的验证负担（第 9 章）。

\textbf{线索三：跨门派密度上升，验证状态断档。}前五行方法的工具几乎全部来自门内；第六行开始，熵、耦合、钉住、Furstenberg 型估计进入；第七行的外部工具密度最高。与之同步的是验证状态从“已证、同行评审”跌到“未形式化、未评审”。两件事未必有因果关系——批量化生成的工作流天然挤占形式化资源——但它们叠在一起，构成了当前的全部风险来源。

\subsection{矩阵之外的一条观察}

矩阵只列了方法自身。如果把“方法之间的可组合性”也当成一个属性，会看到一条容易被忽略的线索：Guth--Wang--Zahl 的流线化（把 Wang--Zahl 体系改写为凸聚类 + 板条语言）让三维结构第一次可以被“打包”成一块引理；OpenAI 三维篇把这块引理再次加权打包；四维篇因此可以直接调用。方法被反复“打包—复用”是工程成熟的标志，这一点上手稿展示了很高的素养（本文判断，\cfM）。代价同样存在：每多一层打包，验证者就多一层黑盒——核验四维篇的最省力路径，是先核验被它当作黑盒的三维引理链（附录 B）。

\begin{keybox}[可记忆主句]
矩阵最后一行的含义是：结论跨度最大、外部工具最密、验证记录最空——三个“最”同时成立，决定了这份工作此刻只能被称为“声称”。
\end{keybox}